\chapter{V4 Paired-Spectrum Orbit Classification}
\label{chap:v4-paired-spectrum-orbits}

\status{Exact finite-orbit classification for the SOH-G014 paired xi spectrum; the exceptional value $F(-1/4)$ remains unresolved.}

\section{Purpose}

SOH-G014 proved that the negative-inversion paired xi spectrum
\[
 P_N:=\{\rho:\xi(\rho)=0,\ \xi(N(\rho))=0\}
\]
is finite.  SOH-G016 then identified its quotient image under
\[
 q(s)=\left(s-\frac12\right)^2
\]
with the paired root set of the quotient entire function $F$.

The present chapter determines the exact orbit structure of $P_N$ under the projective Klein-four algebra already established in Chapters~\ref{chap:euler-riemann-negative-inversion} and~\ref{chap:pauli-spinor-lift}.

No new functional equation is introduced here.  The result is purely a consequence of the existing xi reflection, the exact negative-inversion algebra, and the already proved exclusions at the fixed loci of the other involutions.

\section{The three nontrivial involutions in the $s$-plane}

Let
\[
 R(s)=1-s,
 \qquad
 N(s)=\frac{s-1}{2s-1},
 \qquad
 E(s)=\frac{s}{2s-1}.
\]
These are the $s$-plane representatives of the projective maps
\[
 u\mapsto \frac1u,
 \qquad
 u\mapsto -\frac1u,
 \qquad
 u\mapsto -u
\]
under the Li coordinate $u=s/(1-s)$.

Direct calculation gives
\[
 R^2=N^2=E^2=I,
 \qquad
 RN=NR=E.
\]
Hence
\[
 V_4=\{I,R,N,E\}
\]
acts projectively on the sphere and on the affine chart away from the pole $s=1/2$ of $N$ and $E$.

\section{Invariance of the paired spectrum}

\begin{theorem}[SOH-G017 paired-spectrum invariance]
The finite set $P_N$ is invariant under the full Klein-four action $V_4$.
\end{theorem}

\begin{proof}
Take $s\in P_N$.  Then
\[
 \xi(s)=0,
 \qquad
 \xi(N(s))=0.
\]
The completed xi functional equation gives
\[
 \xi(R(s))=\xi(1-s)=\xi(s)=0.
\]
Because $R$ and $N$ commute,
\[
 N(R(s))=R(N(s)),
\]
and therefore
\[
 \xi(N(R(s)))=\xi(R(N(s)))=\xi(N(s))=0.
\]
Thus $R(s)\in P_N$.

Also, since $N^2=I$,
\[
 \xi(N(s))=0,
 \qquad
 \xi(N(N(s)))=\xi(s)=0,
\]
so $N(s)\in P_N$.  Finally
\[
 E(s)=R(N(s))
\]
and therefore $E(s)\in P_N$ as well.
\end{proof}

\section{Fixed loci}

The possible nongeneric orbit sizes are controlled entirely by stabilizers.  We therefore determine the fixed loci of the three nontrivial involutions.

\subsection{Reflection}

The equation
\[
 R(s)=s
\]
has the unique affine solution
\[
 s=\frac12.
\]
By the positive-kernel result of Chapter~\ref{chap:positive-riemann-kernel},
\[
 \xi\!\left(\frac12\right)
 =\int_0^\infty \Phi(y)\,dy
 >0.
\]
Thus no point of $P_N$ is fixed by $R$.

\subsection{Euler half-turn}

The equation
\[
 E(s)=s
\]
gives
\[
 s\in\{0,1\}.
\]
For the completed xi function,
\[
 \xi(0)=\xi(1)=\frac12.
\]
Hence no point of $P_N$ is fixed by $E$.

\subsection{Negative inversion}

The equation
\[
 N(s)=s
\]
reduces to
\[
 2s^2-2s+1=0,
\]
with solutions
\[
 s_\pm=\frac12\pm\frac{i}{2}.
\]
These two points are exchanged by the reflection:
\[
 R(s_+)=s_-,
 \qquad
 R(s_-)=s_+.
\]
Their xi-zero status is not decided in this chapter.

\section{Complete orbit classification}

\begin{theorem}[SOH-G017 V4 orbit classification]
Every $V_4$-orbit contained in $P_N$ has cardinality $4$, except possibly
\[
 \mathcal O_*
 =\left\{\frac12+\frac i2,\frac12-\frac i2\right\},
\]
which has cardinality $2$ if and only if its points belong to $P_N$.
\end{theorem}

\begin{proof}
For any finite group action,
\[
 |V_4\cdot s|=\frac{|V_4|}{|\operatorname{Stab}(s)|}.
\]
The only possible nontrivial stabilizers are generated by one of the three nonidentity involutions $R$, $N$, or $E$.

The fixed points of $R$ and $E$ are not xi zeros, as shown above.  Therefore a point of $P_N$ can have a nontrivial stabilizer only if it is fixed by $N$.  That forces
\[
 s\in\left\{\frac12+\frac i2,\frac12-\frac i2\right\}.
\]
At either such point the stabilizer has order $2$, generated by $N$, and the orbit therefore has size $4/2=2$.  Every remaining stabilizer is trivial, so every remaining orbit has size $4$.
\end{proof}

\section{Exact cardinality congruence}

Let $a\ge0$ count the generic four-point orbits and let
\[
 \varepsilon\in\{0,1\}
\]
record whether the exceptional two-point orbit is present.  Then
\[
 \boxed{|P_N|=4a+2\varepsilon}.
\]
Equivalently,
\[
 |P_N|\equiv
 \begin{cases}
 0\pmod4,&\varepsilon=0,\\
 2\pmod4,&\varepsilon=1.
 \end{cases}
\]
The exceptional bit satisfies
\[
 \varepsilon=1
 \iff
 \xi\!\left(\frac12+\frac i2\right)=0.
\]

\section{Cross-check with the quotient paired spectrum}

Chapter~\ref{chap:paired-spectrum-quotient-correspondence} proved that
\[
 q(P_N)=P_J
\]
and that the restricted quotient map is exactly two-to-one.  Therefore
\[
 |P_J|=\frac{|P_N|}{2}=2a+\varepsilon.
\]
Thus
\[
 \boxed{|P_J|\equiv\varepsilon\pmod2}.
\]

The exceptional pair maps to
\[
 q\!\left(\frac12\pm\frac i2\right)
 =-\frac14.
\]
Hence
\[
 \varepsilon=1
 \iff
 F(-1/4)=0.
\]
This is a correspondence statement only.  It does not determine whether $F(-1/4)$ vanishes.

\section{What has and has not been proved}

The following statements are now exact:
\begin{itemize}
 \item $P_N$ is a finite $V_4$-invariant set;
 \item every generic paired-spectrum orbit has four points;
 \item the only possible two-point orbit is $\{1/2\pm i/2\}$;
 \item $|P_N|=4a+2\varepsilon$;
 \item $|P_J|=2a+\varepsilon$;
 \item the parity of $P_J$ and the residue of $|P_N|$ modulo four are controlled by the same exceptional bit.
\end{itemize}

The following remain open:
\begin{itemize}
 \item whether $P_N$ or $P_J$ is nonempty;
 \item whether $F(-1/4)=0$;
 \item real-rootedness of $F$;
 \item PF-infinity;
 \item SOH-G003;
 \item the Riemann Hypothesis.
\end{itemize}
